\section{Methodology}
\label{sec:methodology}

This section details our methodology for geolocating cellular infrastructure
using Apple's \ac{CPS}. 

\begin{table}[t]
    \centering
    \caption{OpenCelliD \ac{BTS} counts by cellular generation from March 2025.}
    \label{tab:opencellid-counts}
    \begin{tabular}{c|r}
    \textbf{Technology} & \multicolumn{1}{|c}{\textbf{Count}} \\
    \hline
    GSM & 1,300,596  \\
    UMTS & 473,432 \\
    CDMA & 28 \\
    LTE & 2,726,803 \\
    NR & 120,009  \\
    \hline
    Total & 4,620,869
    \end{tabular}
\end{table}


\subsection{Bootstrapping \acp{BTS}}

Apple's \ac{CPS} requires Apple devices to submit the unique identifiers of
nearby \acp{BTS} in order to receive their geolocations. However, beyond the
\ac{PLMN}, which identifies an \ac{MNO}, there is no universal structure or
format for how cells are numbered. Doing so is left up to individual \acp{MNO},
and their numbering plan may vary substantially from other \acp{MNO}, even
within the same country. Further, because the combined location and cell
identifier address spaces range from 28 to 32 bits, depending on the technology
(\ref{tab:cellids}), guessing within or exhaustively enumerating the identifier
space is an inefficient way to accumulate geolocated cell data.

This presents a chicken and egg problem---how might one gather a baseline of
geolocated \ac{BTS} data to bootstrap the geolocation? To solve this dilemma, we
turned to a free database provided by OpenCelliD~\cite{opencellid}. OpenCelliD
provides worldwide, geolocated \ac{BTS} data for free
